% Native editable Beamer slides. Build from the repository root.
\documentclass[aspectratio=169,11pt,handout,t]{beamer}
\makeatletter
\def\input@path{{../}}
\makeatother
\usepackage{satnetslides}
\usepackage{listings}
\definecolor{Icon0}{HTML}{555555}
\definecolor{Icon1}{HTML}{999999}
\definecolor{Icon2}{HTML}{111111}
\definecolor{Icon3}{HTML}{F3F3F3}
\definecolor{Icon4}{HTML}{BDBDBD}
\definecolor{Icon5}{HTML}{666666}
\definecolor{Icon6}{HTML}{F8F8F8}
\definecolor{Icon7}{HTML}{B8B8B8}
\definecolor{Icon8}{HTML}{C5C5C5}
\definecolor{Icon9}{HTML}{DEDEDE}
\definecolor{Icon10}{HTML}{A8A8A8}
\definecolor{Icon11}{HTML}{FAFAFA}
\definecolor{Icon12}{HTML}{8E8E8E}
\definecolor{Icon13}{HTML}{E6E6E6}
\definecolor{Icon14}{HTML}{CCCCCC}
\definecolor{Icon15}{HTML}{E2E2E2}
\definecolor{Icon16}{HTML}{D5D5D5}
\definecolor{Icon17}{HTML}{F7F7F7}
\definecolor{Icon18}{HTML}{DADADA}
\definecolor{Icon19}{HTML}{777777}
\definecolor{Icon20}{HTML}{DDDDDD}
\definecolor{Icon21}{HTML}{C8C8C8}
\definecolor{Icon22}{HTML}{F0F0F0}
\definecolor{Icon23}{HTML}{CFCFCF}
\definecolor{Icon24}{HTML}{C6C6C6}
\definecolor{Icon25}{HTML}{E0E0E0}
\definecolor{Icon26}{HTML}{E5E5E5}
\definecolor{Icon27}{HTML}{4B4B4B}
\definecolor{Icon28}{HTML}{FFFFFF}
\definecolor{Icon29}{HTML}{BEBEBE}
\definecolor{Icon30}{HTML}{F4F4F4}
\definecolor{Icon31}{HTML}{ECECEC}
% Native TikZ paths for the shared monochrome component icons.
% Both solar arrays use the same axonometric projection.
\tikzset{pics/satellite/.style={code={%
\path[fill=none,draw=Icon0,line width=1.14pt,line join=round,line cap=round] (-0.276,-0.00092) -- (0.276,0.12052);
\path[fill=Icon1,draw=Icon2,line width=0.304pt,line join=round,line cap=round] (-0.5244,-0.15346) -- (-0.184,-0.07857) -- (-0.184,-0.09697) -- (-0.5244,-0.17186) -- cycle;
\path[fill=Icon3,draw=Icon2,line width=0.494pt,line join=round,line cap=round] (-0.6532,0.01398) -- (-0.3128,0.08887) -- (-0.184,-0.07857) -- (-0.5244,-0.15346) -- cycle;
\path[fill=Icon4,draw=Icon5,line width=0.19pt,line join=round,line cap=round] (-0.62836,0.00546) -- (-0.31924,0.07347) -- (-0.20884,-0.07005) -- (-0.51796,-0.13806) -- cycle;
\path[fill=none,draw=Icon6,line width=0.247pt,line join=round,line cap=round] (-0.57684,0.0168) -- (-0.46644,-0.12672);
\path[fill=none,draw=Icon6,line width=0.247pt,line join=round,line cap=round] (-0.52532,0.02813) -- (-0.41492,-0.11539);
\path[fill=none,draw=Icon6,line width=0.247pt,line join=round,line cap=round] (-0.4738,0.03947) -- (-0.3634,-0.10405);
\path[fill=none,draw=Icon6,line width=0.247pt,line join=round,line cap=round] (-0.42228,0.0508) -- (-0.31188,-0.09272);
\path[fill=none,draw=Icon6,line width=0.247pt,line join=round,line cap=round] (-0.37076,0.06214) -- (-0.26036,-0.08138);
\path[fill=none,draw=Icon6,line width=0.247pt,line join=round,line cap=round] (-0.59156,-0.04238) -- (-0.28244,0.02563);
\path[fill=none,draw=Icon6,line width=0.247pt,line join=round,line cap=round] (-0.55476,-0.09022) -- (-0.24564,-0.02221);
\path[fill=none,draw=Icon5,line width=0.304pt,line join=round,line cap=round] (-0.483,0.05143) -- (-0.3542,-0.11601);
\path[fill=Icon1,draw=Icon2,line width=0.304pt,line join=round,line cap=round] (0.3128,0.03073) -- (0.6532,0.10562) -- (0.6532,0.08722) -- (0.3128,0.01233) -- cycle;
\path[fill=Icon3,draw=Icon2,line width=0.494pt,line join=round,line cap=round] (0.184,0.19817) -- (0.5244,0.27306) -- (0.6532,0.10562) -- (0.3128,0.03073) -- cycle;
\path[fill=Icon4,draw=Icon5,line width=0.19pt,line join=round,line cap=round] (0.20884,0.18965) -- (0.51796,0.25766) -- (0.62836,0.11414) -- (0.31924,0.04613) -- cycle;
\path[fill=none,draw=Icon6,line width=0.247pt,line join=round,line cap=round] (0.26036,0.20098) -- (0.37076,0.05746);
\path[fill=none,draw=Icon6,line width=0.247pt,line join=round,line cap=round] (0.31188,0.21232) -- (0.42228,0.0688);
\path[fill=none,draw=Icon6,line width=0.247pt,line join=round,line cap=round] (0.3634,0.22365) -- (0.4738,0.08013);
\path[fill=none,draw=Icon6,line width=0.247pt,line join=round,line cap=round] (0.41492,0.23499) -- (0.52532,0.09147);
\path[fill=none,draw=Icon6,line width=0.247pt,line join=round,line cap=round] (0.46644,0.24632) -- (0.57684,0.1028);
\path[fill=none,draw=Icon6,line width=0.247pt,line join=round,line cap=round] (0.24564,0.14181) -- (0.55476,0.20982);
\path[fill=none,draw=Icon6,line width=0.247pt,line join=round,line cap=round] (0.28244,0.09397) -- (0.59156,0.16198);
\path[fill=none,draw=Icon5,line width=0.304pt,line join=round,line cap=round] (0.3542,0.23561) -- (0.483,0.06817);
\path[fill=Icon7,draw=Icon2,line width=0.418pt,line join=round,line cap=round] (-0.1978,0.18961) -- (-0.115,0.08197) -- (-0.115,-0.07443) -- (-0.1978,0.03321) -- cycle;
\path[fill=Icon8,draw=Icon2,line width=0.418pt,line join=round,line cap=round] (-0.115,0.08197) -- (-0.0736,0.0701) -- (-0.0736,-0.0863) -- (-0.115,-0.07443) -- cycle;
\path[fill=Icon9,draw=Icon2,line width=0.418pt,line join=round,line cap=round] (-0.0736,0.0701) -- (0.184,0.12678) -- (0.184,-0.02962) -- (-0.0736,-0.0863) -- cycle;
\path[fill=Icon7,draw=Icon2,line width=0.418pt,line join=round,line cap=round] (0.184,0.12678) -- (0.1978,0.15079) -- (0.1978,-0.00561) -- (0.184,-0.02962) -- cycle;
\path[fill=Icon10,draw=Icon2,line width=0.418pt,line join=round,line cap=round] (0.1978,0.15079) -- (0.115,0.25843) -- (0.115,0.10203) -- (0.1978,-0.00561) -- cycle;
\path[fill=Icon11,draw=Icon2,line width=0.475pt,line join=round,line cap=round] (-0.184,0.21362) -- (0.0736,0.2703) -- (0.115,0.25843) -- (0.1978,0.15079) -- (0.184,0.12678) -- (-0.0736,0.0701) -- (-0.115,0.08197) -- (-0.1978,0.18961) -- cycle;
\path[fill=Icon12,draw=Icon2,line width=0.266pt,line join=round,line cap=round] (-0.0276,0.04342) -- (0.1472,0.08188) -- (0.1472,-0.00092) -- (-0.0276,-0.03938) -- cycle;
\path[fill=none,draw=Icon13,line width=0.209pt,line join=round,line cap=round] (0.0161,0.04384) -- (0.0161,-0.02056);
\path[fill=none,draw=Icon13,line width=0.209pt,line join=round,line cap=round] (0.0598,0.05345) -- (0.0598,-0.01095);
\path[fill=none,draw=Icon13,line width=0.209pt,line join=round,line cap=round] (0.1035,0.06307) -- (0.1035,-0.00133);
\path[fill=Icon14,draw=Icon2,line width=0.304pt,line join=round,line cap=round] (0.0782,0.20599) -- (0.08144,0.20056) -- (0.08297,0.195) -- (0.08272,0.1895) -- (0.08071,0.18429) -- (0.07701,0.17957) -- (0.07177,0.1755) -- (0.06518,0.17226) -- (0.0575,0.16997) -- (0.04903,0.16871) -- (0.04008,0.16853) -- (0.03101,0.16945) -- (0.02216,0.17141) -- (0.01388,0.17436) -- (0.00647,0.17817) -- (0.00023,0.1827) -- (-0.0046,0.18777) -- (-0.00784,0.1932) -- (-0.00937,0.19876) -- (-0.00912,0.20426) -- (-0.00711,0.20947) -- (-0.00341,0.21419) -- (0.00183,0.21826) -- (0.00842,0.2215) -- (0.0161,0.22379) -- (0.02457,0.22505) -- (0.03352,0.22523) -- (0.04259,0.22431) -- (0.05144,0.22235) -- (0.05972,0.2194) -- (0.06713,0.21559) -- (0.07337,0.21106) -- cycle;
\path[fill=none,draw=Icon2,line width=0.437pt,line join=round,line cap=round] (-0.1012,0.1759) -- (-0.1012,0.2955);
\path[fill=none,draw=Icon2,line width=0.38pt,line join=round,line cap=round] (-0.1288,0.28943) -- (-0.0736,0.30158);
}}}
\tikzset{pics/user/.style={code={%
\path[fill=Icon15,draw=Icon2,line width=0.494pt,line join=round,line cap=round] (-0.2408,0.172) ellipse [x radius=0.0688cm,y radius=0.0688cm];
\path[fill=Icon16,draw=Icon2,line width=0.532pt,line join=round,line cap=round] (-0.3612,-0.0946) -- (-0.3526,0.0086) -- (-0.3096,0.0602) -- (-0.1806,0.0602) -- (-0.1204,0.0086) -- (-0.086,-0.0946) -- cycle;
\path[fill=Icon17,draw=Icon2,line width=0.532pt,line join=round,line cap=round] (-0.1118,0.0946) -- (0.215,0.0946) -- (0.1806,-0.1376) -- (-0.1376,-0.1376) -- cycle;
\path[fill=Icon18,draw=Icon19,line width=0.304pt,line join=round,line cap=round] (-0.0774,0.0602) -- (0.172,0.0602) -- (0.1462,-0.0946) -- (-0.1032,-0.0946) -- cycle;
\path[fill=Icon20,draw=Icon2,line width=0.532pt,line join=round,line cap=round] (-0.1376,-0.1376) -- (0.1806,-0.1376) -- (0.2752,-0.1978) -- (-0.2408,-0.1978) -- cycle;
\path[fill=none,draw=Icon19,line width=0.323pt,line join=round,line cap=round] (-0.1032,-0.1634) -- (0.1032,-0.1634);
\path[fill=none,draw=Icon2,line width=0.57pt,line join=round,line cap=round] (-0.3698,-0.2322) -- (0.3096,-0.2322);
}}}
\tikzset{pics/user-router/.style={code={%
\path[fill=Icon15,draw=Icon2,line width=0.494pt,line join=round,line cap=round] (-0.2408,0.172) ellipse [x radius=0.0688cm,y radius=0.0688cm];
\path[fill=Icon16,draw=Icon2,line width=0.532pt,line join=round,line cap=round] (-0.3612,-0.0946) -- (-0.3526,0.0086) -- (-0.3096,0.0602) -- (-0.1806,0.0602) -- (-0.1204,0.0086) -- (-0.086,-0.0946) -- cycle;
\path[fill=Icon17,draw=Icon2,line width=0.532pt,line join=round,line cap=round] (-0.1118,0.0946) -- (0.215,0.0946) -- (0.1806,-0.1376) -- (-0.1376,-0.1376) -- cycle;
\path[fill=Icon18,draw=Icon19,line width=0.304pt,line join=round,line cap=round] (-0.0774,0.0602) -- (0.172,0.0602) -- (0.1462,-0.0946) -- (-0.1032,-0.0946) -- cycle;
\path[fill=Icon20,draw=Icon2,line width=0.532pt,line join=round,line cap=round] (-0.1376,-0.1376) -- (0.1806,-0.1376) -- (0.2752,-0.1978) -- (-0.2408,-0.1978) -- cycle;
\path[fill=none,draw=Icon19,line width=0.323pt,line join=round,line cap=round] (-0.1032,-0.1634) -- (0.1032,-0.1634);
\path[fill=none,draw=Icon2,line width=0.57pt,line join=round,line cap=round] (-0.3698,-0.2322) -- (0.3096,-0.2322);
\path[fill=Icon21,draw=Icon2,line width=0.494pt,line join=round,line cap=round] (0.2322,-0.0258) -- (0.3956,-0.0258) -- (0.3956,-0.1376) -- (0.2322,-0.1376) -- cycle;
\path[fill=none,draw=Icon2,line width=0.57pt,line join=round,line cap=round] (0.344,-0.0258) -- (0.344,0.1032);
\path[fill=Icon2,draw=Icon2,line width=0.152pt,line join=round,line cap=round] (0.27348,-0.08428) ellipse [x radius=0.00688cm,y radius=0.00688cm];
\path[fill=Icon2,draw=Icon2,line width=0.152pt,line join=round,line cap=round] (0.31648,-0.08428) ellipse [x radius=0.00688cm,y radius=0.00688cm];
\path[fill=Icon2,draw=Icon2,line width=0.152pt,line join=round,line cap=round] (0.35948,-0.08428) ellipse [x radius=0.00688cm,y radius=0.00688cm];
}}}
\tikzset{pics/terminal/.style={code={%
\path[fill=Icon4,draw=Icon2,line width=0.532pt,line join=round,line cap=round] (-0.215,0.1892) -- (0.258,0.1204) -- (0.215,-0.043) -- (-0.258,0.0258) -- cycle;
\path[fill=Icon22,draw=Icon2,line width=0.532pt,line join=round,line cap=round] (-0.215,0.215) -- (0.258,0.1462) -- (0.2236,-0.0086) -- (-0.2494,0.0602) -- cycle;
\path[fill=none,draw=Icon1,line width=0.285pt,line join=round,line cap=round] (-0.1806,0.1806) -- (0.215,0.1204);
\path[fill=none,draw=Icon0,line width=1.14pt,line join=round,line cap=round] (-0.0086,0.0172) -- (-0.0086,-0.1376);
\path[fill=Icon23,draw=Icon2,line width=0.532pt,line join=round,line cap=round] (-0.0344,-0.1204) -- (0.0172,-0.1204) -- (0.1204,-0.2408) -- (0.0516,-0.2408) -- (-0.0086,-0.1634) -- (-0.1376,-0.2408) -- (-0.215,-0.2408) -- cycle;
\path[fill=none,draw=Icon2,line width=0.76pt,line join=round,line cap=round] (-0.0086,-0.1376) -- (0.1548,-0.1806);
}}}
\tikzset{pics/gateway/.style={code={%
\path[fill=Icon24,draw=Icon2,line width=0.532pt,line join=round,line cap=round] (-0.1462,-0.1978) -- (0.1376,-0.1978) -- (0.1978,-0.258) -- (-0.215,-0.258) -- cycle;
\path[fill=Icon16,draw=Icon2,line width=0.532pt,line join=round,line cap=round] (-0.0258,-0.0258) -- (0.0516,-0.0344) -- (0.0946,-0.1978) -- (-0.086,-0.1978) -- cycle;
\path[fill=Icon25,draw=Icon2,line width=0.532pt,line join=round,line cap=round] (-0.25717,0.03102) -- (-0.23001,0.00839) -- (-0.20385,-0.0121) -- (-0.1787,-0.03043) -- (-0.15454,-0.04662) -- (-0.13139,-0.06066) -- (-0.10924,-0.07255) -- (-0.08809,-0.08229) -- (-0.06794,-0.08989) -- (-0.04879,-0.09534) -- (-0.03065,-0.09865) -- (-0.0135,-0.0998) -- (0.00264,-0.09881) -- (0.01778,-0.09567) -- (0.03192,-0.09038) -- (0.04506,-0.08295) -- (0.0572,-0.07337) -- (0.06834,-0.06164) -- (0.07847,-0.04776) -- (0.08761,-0.03174) -- (0.09574,-0.01357) -- (0.10287,0.00675) -- (0.109,0.02922) -- (0.11413,0.05383) -- (0.11825,0.08059) -- (0.12138,0.1095) -- (0.12351,0.14055) -- (0.12463,0.17376) -- (0.12475,0.20911) -- cycle;
\path[fill=Icon11,draw=Icon2,line width=0.532pt,line join=round,line cap=round] (0.12475,0.20911) -- (0.12581,0.2007) -- (0.12214,0.1903) -- (0.11384,0.17817) -- (0.1011,0.16461) -- (0.08424,0.14996) -- (0.06368,0.13457) -- (0.03991,0.11882) -- (0.01354,0.1031) -- (-0.0148,0.0878) -- (-0.0444,0.0733) -- (-0.07454,0.05994) -- (-0.10448,0.04807) -- (-0.13347,0.03797) -- (-0.16081,0.02989) -- (-0.18582,0.02403) -- (-0.20788,0.02053) -- (-0.22645,0.01949) -- (-0.24108,0.02092) -- (-0.25141,0.0248) -- (-0.25717,0.03102) -- (-0.25823,0.03943) -- (-0.25456,0.04983) -- (-0.24626,0.06195) -- (-0.23352,0.07551) -- (-0.21666,0.09017) -- (-0.19609,0.10556) -- (-0.17233,0.1213) -- (-0.14596,0.13702) -- (-0.11762,0.15232) -- (-0.08802,0.16683) -- (-0.05788,0.18018) -- (-0.02794,0.19206) -- (0.00105,0.20216) -- (0.02839,0.21024) -- (0.0534,0.21609) -- (0.07546,0.21959) -- (0.09404,0.22063) -- (0.10866,0.2192) -- (0.11899,0.21533) -- cycle;
\path[fill=none,draw=Icon2,line width=0.437pt,line join=round,line cap=round] (-0.24158,0.03829) -- (-0.11346,0.22139);
\path[fill=none,draw=Icon2,line width=0.437pt,line join=round,line cap=round] (0.10916,0.20184) -- (-0.11346,0.22139);
\path[fill=Icon5,draw=Icon2,line width=0.494pt,line join=round,line cap=round] (-0.11346,0.22139) ellipse [x radius=0.0215cm,y radius=0.0215cm];
\path[fill=Icon26,draw=Icon2,line width=0.494pt,line join=round,line cap=round] (0.2322,-0.086) -- (0.3784,-0.086) -- (0.3784,-0.2494) -- (0.2322,-0.2494) -- cycle;
\path[fill=none,draw=Icon19,line width=0.323pt,line join=round,line cap=round] (0.258,-0.129) -- (0.344,-0.129);
\path[fill=none,draw=Icon19,line width=0.323pt,line join=round,line cap=round] (0.258,-0.1634) -- (0.344,-0.1634);
\path[fill=none,draw=Icon19,line width=0.323pt,line join=round,line cap=round] (0.258,-0.1978) -- (0.344,-0.1978);
}}}
\tikzset{pics/internet/.style={code={%
\path[fill=Icon17,draw=Icon2,line width=0.494pt,line join=round,line cap=round] (0,0.0086) ellipse [x radius=0.2408cm,y radius=0.2408cm];
\path[fill=none,draw=Icon19,line width=0.38pt,line join=round,line cap=round] (0,0.0086) ellipse [x radius=0.129cm,y radius=0.2408cm];
\path[fill=none,draw=Icon19,line width=0.38pt,line join=round,line cap=round] (0,0.0086) ellipse [x radius=0.2408cm,y radius=0.0946cm];
\path[fill=none,draw=Icon1,line width=0.342pt,line join=round,line cap=round] (-0.2408,0.0086) -- (0.2408,0.0086);
\path[fill=none,draw=Icon1,line width=0.342pt,line join=round,line cap=round] (0,0.2494) -- (0,-0.2322);
\path[fill=none,draw=Icon2,line width=0.475pt,line join=round,line cap=round] (-0.1806,0.1204) -- (0.1032,0.172);
\path[fill=none,draw=Icon2,line width=0.475pt,line join=round,line cap=round] (0.1032,0.172) -- (0.172,-0.0946);
\path[fill=none,draw=Icon2,line width=0.475pt,line join=round,line cap=round] (0.172,-0.0946) -- (-0.1204,-0.1462);
\path[fill=none,draw=Icon2,line width=0.475pt,line join=round,line cap=round] (-0.1204,-0.1462) -- (-0.1806,0.1204);
\path[fill=Icon27,draw=Icon28,line width=0.304pt,line join=round,line cap=round] (-0.1806,0.1204) ellipse [x radius=0.0258cm,y radius=0.0258cm];
\path[fill=Icon27,draw=Icon28,line width=0.304pt,line join=round,line cap=round] (0.1032,0.172) ellipse [x radius=0.0258cm,y radius=0.0258cm];
\path[fill=Icon27,draw=Icon28,line width=0.304pt,line join=round,line cap=round] (0.172,-0.0946) ellipse [x radius=0.0258cm,y radius=0.0258cm];
\path[fill=Icon27,draw=Icon28,line width=0.304pt,line join=round,line cap=round] (-0.1204,-0.1462) ellipse [x radius=0.0258cm,y radius=0.0258cm];
}}}
\tikzset{pics/server/.style={code={%
\path[fill=Icon29,draw=Icon2,line width=0.532pt,line join=round,line cap=round] (0.1462,0.2064) -- (0.2494,0.1376) -- (0.2494,-0.258) -- (0.1462,-0.1978) -- cycle;
\path[fill=Icon30,draw=Icon2,line width=0.532pt,line join=round,line cap=round] (-0.2064,0.2064) -- (0.1462,0.2064) -- (0.2494,0.1376) -- (-0.1032,0.1376) -- cycle;
\path[fill=Icon31,draw=Icon2,line width=0.494pt,line join=round,line cap=round] (-0.2064,0.2064) -- (0.1462,0.2064) -- (0.1462,-0.215) -- (-0.2064,-0.215) -- cycle;
\path[fill=Icon28,draw=Icon0,line width=0.38pt,line join=round,line cap=round] (-0.172,0.1548) -- (0.1118,0.1548) -- (0.1118,0.0688) -- (-0.172,0.0688) -- cycle;
\path[fill=Icon5,draw=none,line width=0pt,line join=round,line cap=round] (-0.13072,0.1118) ellipse [x radius=0.01548cm,y radius=0.01548cm];
\path[fill=none,draw=Icon19,line width=0.304pt,line join=round,line cap=round] (-0.0602,0.129) -- (-0.0602,0.0946);
\path[fill=none,draw=Icon19,line width=0.304pt,line join=round,line cap=round] (-0.0258,0.129) -- (-0.0258,0.0946);
\path[fill=none,draw=Icon19,line width=0.304pt,line join=round,line cap=round] (0.0086,0.129) -- (0.0086,0.0946);
\path[fill=none,draw=Icon19,line width=0.304pt,line join=round,line cap=round] (0.043,0.129) -- (0.043,0.0946);
\path[fill=Icon28,draw=Icon0,line width=0.38pt,line join=round,line cap=round] (-0.172,0.043) -- (0.1118,0.043) -- (0.1118,-0.043) -- (-0.172,-0.043) -- cycle;
\path[fill=Icon5,draw=none,line width=0pt,line join=round,line cap=round] (-0.13072,0) ellipse [x radius=0.01548cm,y radius=0.01548cm];
\path[fill=none,draw=Icon19,line width=0.304pt,line join=round,line cap=round] (-0.0602,0.0172) -- (-0.0602,-0.0172);
\path[fill=none,draw=Icon19,line width=0.304pt,line join=round,line cap=round] (-0.0258,0.0172) -- (-0.0258,-0.0172);
\path[fill=none,draw=Icon19,line width=0.304pt,line join=round,line cap=round] (0.0086,0.0172) -- (0.0086,-0.0172);
\path[fill=none,draw=Icon19,line width=0.304pt,line join=round,line cap=round] (0.043,0.0172) -- (0.043,-0.0172);
\path[fill=Icon28,draw=Icon0,line width=0.38pt,line join=round,line cap=round] (-0.172,-0.0688) -- (0.1118,-0.0688) -- (0.1118,-0.1548) -- (-0.172,-0.1548) -- cycle;
\path[fill=Icon5,draw=none,line width=0pt,line join=round,line cap=round] (-0.13072,-0.1118) ellipse [x radius=0.01548cm,y radius=0.01548cm];
\path[fill=none,draw=Icon19,line width=0.304pt,line join=round,line cap=round] (-0.0602,-0.0946) -- (-0.0602,-0.129);
\path[fill=none,draw=Icon19,line width=0.304pt,line join=round,line cap=round] (-0.0258,-0.0946) -- (-0.0258,-0.129);
\path[fill=none,draw=Icon19,line width=0.304pt,line join=round,line cap=round] (0.0086,-0.0946) -- (0.0086,-0.129);
\path[fill=none,draw=Icon19,line width=0.304pt,line join=round,line cap=round] (0.043,-0.0946) -- (0.043,-0.129);
\path[fill=none,draw=Icon2,line width=0.874pt,line join=round,line cap=round] (-0.1806,-0.2322) -- (-0.1032,-0.2322);
\path[fill=none,draw=Icon2,line width=0.874pt,line join=round,line cap=round] (0.0516,-0.2322) -- (0.129,-0.2322);
}}}

\lstset{language=Python,basicstyle=\ttfamily\fontsize{7.5}{9.5}\selectfont,
  keywordstyle=\bfseries\color{Ink},commentstyle=\color{Muted},
  stringstyle=\color{Muted},showstringspaces=false,columns=fullflexible,
  keepspaces=true,frame=single,rulecolor=\color{Line},xleftmargin=3pt,xrightmargin=3pt}
\tikzset{box/.style={draw=Line,fill=Soft,rounded corners=1pt,font=\scriptsize,
  align=center,minimum height=0.65cm,inner sep=4pt},
  dot/.style={circle,draw=Ink,fill=Paper,minimum size=0.50cm,inner sep=0pt,font=\scriptsize\bfseries}}
\pgfplotsset{courseaxis/.style={width=6.6cm,height=4.0cm,
  label style={font=\scriptsize},tick label style={font=\scriptsize},
  axis line style={Muted},tick style={Muted},grid=major,grid style={Line,line width=0.25pt},
  legend style={font=\scriptsize,draw=none,fill=Paper},unbounded coords=jump}}
\setlecture{04}{Dynamics, Handover, and Performance}
\title{Dynamics, Handover, and Performance}
\author{Yi Ching (David) Chou}
\date{}
\begin{document}
\begin{frame}[plain]
  \begin{tikzpicture}[remember picture,overlay]
    \fill[Paper] (current page.south west) rectangle (current page.north east);
    \fill[Line] (current page.south west) rectangle ([yshift=0.18cm]current page.south east);
    \fill[Ink] ([xshift=0.95cm,yshift=-0.95cm]current page.north west) rectangle ++(0.12cm,-5.35cm);
    \node[anchor=north west,text=Muted] at ([xshift=1.35cm,yshift=-0.9cm]current page.north west)
      {\small\bfseries SATNET EDU / Lecture 04};
    \node[anchor=north west,align=left,text width=12.2cm]
      at ([xshift=1.35cm,yshift=-1.65cm]current page.north west)
      {{\fontsize{29}{33}\selectfont\bfseries Dynamics, Handover,\\[0.12cm]and Performance}};
    \node[anchor=north west,align=left,text width=11.6cm,text=Muted]
      at ([xshift=1.38cm,yshift=-3.85cm]current page.north west)
      {\small How changing links and routes affect communication over time.};
    \node[anchor=north west,fill=Ink,rounded corners=2pt,inner xsep=0.18cm,inner ysep=0.09cm,text=Paper]
      (tag) at ([xshift=1.38cm,yshift=-4.78cm]current page.north west)
      {\small\bfseries LECTURE 04};
    \node[anchor=west,align=left,text width=9.5cm] at ([xshift=0.28cm]tag.east)
      {\small From a working snapshot to a changing connection};
    \draw[Ink,line width=1.1pt] ([xshift=1.38cm,yshift=-5.55cm]current page.north west) -- ++(10.9cm,0);
    \node[anchor=north,align=center,text=Muted] at ([yshift=-5.85cm]current page.north)
      {\small Follow the timeline. Explain the transition. Measure the effect.};
    \node[anchor=north,align=center] at ([yshift=-6.65cm]current page.north)
      {\small Instructor: Yi Ching (David) Chou};
  \end{tikzpicture}
\end{frame}

\begin{frame}[fragile]{A Working Route Is a Statement About One Time}
  \context{Satellite motion changes link availability and distance even when both ground terminals stay still.}
  \begin{columns}[T,onlytextwidth]
    \begin{column}{0.47\textwidth}\raggedright
      \begin{teachingpoints}
        \point{Edges can appear or disappear.}{\item A terminal crosses a footprint edge. An ISL exceeds its range. A scheduled failure disables a node.}
        \point{Weights can change without an outage.}{\item The same sequence of satellites can remain connected while changing distances alter its propagation delay.}
        \point{Track the graph through time.}{\item A route at time $t$ does not prove that the same path will work at $t+\Delta t$.}
      \end{teachingpoints}
    \end{column}
    \begin{column}{0.49\textwidth}\centering
      \figtitle{Same endpoints, a different usable chain}
\begin{tikzpicture}[x=0.90cm,y=0.90cm]
\path[use as bounding box] (-0.25,-0.2) rectangle (7.15,4.5);

\foreach \yy/\id/\lab in {3.5/first/{Time $t_0$},1.2/second/{Time $t_1$}}{
 \node[smalllab,anchor=west] at (0,\yy+0.65) {\lab};
 \node[dot] (u\id) at (0.55,\yy) {U};
 \node[dot] (a\id) at (2.3,\yy) {A};
 \node[dot] (b\id) at (4.1,\yy) {B};
 \node[dot] (v\id) at (5.95,\yy) {V};
}
\draw[net,line width=1pt] (ufirst)--(afirst)--(bfirst)--(vfirst);
\draw[net,dashed,Muted] (usecond)--(asecond);
\draw[net,line width=1pt] (usecond) to[bend right=25] (bsecond);
\draw[net] (asecond)--(bsecond)--(vsecond);
\node[smalllab] at (3.15,0.05) {U now enters through B};

\end{tikzpicture}
    \end{column}
  \end{columns}
  \assumption{Constructed snapshots. Dashed lines mark an unavailable former link. No control-plane timing is implied.}
  \note{A snapshot route is an instantaneous graph answer. It does not promise that all edges remain available during a packet journey or model routing convergence.}
\end{frame}

\begin{frame}[fragile]{Handover Transfers Service Between Access Links}
  \context{Overlapping contact windows provide time to prepare a new access link before the old one ends.}
  \begin{columns}[T,onlytextwidth]
    \begin{column}{0.47\textwidth}\raggedright
      \begin{teachingpoints}
        \point{Separate availability from selection.}{\item A terminal may see both A and B, while its traffic uses only A. Visibility does not specify a handover policy.}
        \point{Use the overlap to prepare.}{\item Acquire the new link, validate it, and transfer traffic before releasing the old link when hardware permits.}
        \point{Account for setup time.}{\item A 20-second overlap with 6 seconds of preparation leaves at most 14 seconds for other actions and margin.}
      \end{teachingpoints}
    \end{column}
    \begin{column}{0.49\textwidth}\centering
      \figtitle{An overlap creates a preparation opportunity}
\begin{tikzpicture}[x=0.90cm,y=0.90cm]
\path[use as bounding box] (-0.25,-0.2) rectangle (7.15,4.5);

\foreach \yy/\name in {3.7/{A visible},2.8/{B visible},1.8/{Prepare B},0.85/{Selected}}{
 \node[smalllab,anchor=east] at (1.35,\yy) {\name};
 \draw[faint] (1.5,\yy)--(6.8,\yy);
}
\fill[Soft] (3.5,0.45) rectangle (5.5,4.05);
\draw[Ink,line width=3pt] (1.5,3.7)--(5.5,3.7);
\draw[Muted,line width=3pt] (3.5,2.8)--(6.8,2.8);
\draw[Ink,line width=3pt] (3.5,1.8)--(4.1,1.8);
\node[smalllab,above] at (3.8,1.8) {6 s};
\draw[Ink,line width=3pt] (1.5,0.85)--(4.1,0.85);
\draw[Muted,line width=3pt] (4.1,0.85)--(6.8,0.85);
\node[smalllab,above] at (2.4,0.85) {A};\node[smalllab,above] at (5.7,0.85) {B};
\foreach \x/\t in {1.5/0,3.5/20,4.1/26,5.5/40,6.8/53}{\node[smalllab,below] at (\x,0.25) {\t};}
\node[smalllab] at (4.5,4.3) {20 s overlap};
\node[smalllab,anchor=east] at (1.35,0) {Time (s)};

\end{tikzpicture}
    \end{column}
  \end{columns}
  \assumption{Constructed make-before-break example. Dual-link preparation requires supporting terminal hardware.}
  \note{A is available from 0 to 40 seconds and B from 20 to 53. Preparation occupies 20 to 26. Switching at 26 occurs before A leaves. Current SatNet Edu records available links, not this handover state machine.}
\end{frame}

\begin{frame}[fragile]{The Best Link Now May Leave Soon}
  \context{An access policy can consider both current geometry and the time remaining before contact ends.}
  \begin{columns}[T,onlytextwidth]
    \begin{column}{0.47\textwidth}\raggedright
      \begin{teachingpoints}
        \point{High elevation can shorten the link.}{\item It often gives a shorter slant range. A descending satellite can still be close to the end of its contact.}
        \point{Remaining contact time is different.}{\item A rising satellite may have a longer usable future window even when its current elevation is lower.}
        \point{Avoid switching on every small gain.}{\item A hysteresis margin requires a meaningful improvement. A minimum dwell time can also reduce repeated switching.}
      \end{teachingpoints}
    \end{column}
    \begin{column}{0.49\textwidth}\centering
      \figtitle{Two candidates at the same decision time}
{\scriptsize\setlength{\tabcolsep}{4pt}\renewcommand{\arraystretch}{1.35}
\begin{tabular}{@{}lll@{}}\toprule
\textbf{Candidate} & \textbf{Elevation} & \textbf{Time left} \\ \midrule
A, descending & $55^\circ$ & 35 s \\
B, rising & $40^\circ$ & 180 s \\
\bottomrule\end{tabular}\par}
\par\vspace{0.3cm}
\begin{tikzpicture}[x=0.9cm,y=0.9cm]
\node[box,text width=2.4cm] at (1.55,1.55) {Best elevation now\\chooses A};
\node[box,text width=2.4cm] at (5.0,1.55) {Longer contact\\chooses B};
\node[lab,text width=5.7cm] at (3.3,0.1) {Example rule: switch only if the new score exceeds the old score by margin $H$};
\end{tikzpicture}

    \end{column}
  \end{columns}
  \assumption{Illustrative candidate values and policies. Future visibility does not guarantee capacity or route availability.}
  \note{The score must be defined in a real policy and may combine elevation, predicted contact duration, load, and route quality. Hysteresis and minimum dwell time trade responsiveness against switching. No particular operator policy is claimed.}
\end{frame}

\begin{frame}[fragile]{Path Changes Can Make Delay Jump}
  \context{Orbital motion changes a path smoothly, but selecting a different path can cause a sudden shift.}
  \begin{columns}[T,onlytextwidth]
    \begin{column}{0.47\textwidth}\raggedright
      \begin{teachingpoints}
        \point{Follow each candidate path.}{\item Its propagation cost varies as the satellites move. Availability can end before its cost becomes unattractive.}
        \point{A switch can add delay.}{\item Here the 12 ms path disappears at 30 seconds. The available replacement costs 17 ms at that moment.}
        \point{Report changes as well as averages.}{\item Two traces with the same mean delay can have different jumps, interruptions, and route lifetimes.}
      \end{teachingpoints}
    \end{column}
    \begin{column}{0.49\textwidth}\centering
      \figtitle{The cheapest path becomes unavailable}

\begin{tikzpicture}
\begin{axis}[courseaxis,xmin=0,xmax=60,ymin=8,ymax=22,
xtick={0,15,30,45,60},ytick={10,12,15,17,20},
xlabel={Time (s)},ylabel={Propagation delay (ms)}]
\addplot[Muted,dashed,line width=0.8pt] coordinates {(0,19)(30,17)(60,15)};
\addplot[Ink,line width=1.3pt] coordinates {(0,10)(30,12)(30,17)(60,15)};
\node[font=\scriptsize,anchor=west] at (axis cs:3,11.5) {Path A};
\node[font=\scriptsize,anchor=west] at (axis cs:38,18.5) {Path B};
\draw[Muted,densely dotted] (axis cs:30,8)--(axis cs:30,22);
\end{axis}
\end{tikzpicture}
\par\vspace{0.12cm}{\footnotesize Bold: selected path\quad Dashed: available alternative\par}
    \end{column}
  \end{columns}
  \assumption{Constructed path costs. The jump is 5 ms. No packet loss or queueing is inferred from this plot.}
  \note{At 30 seconds path A ceases to be available. The bold vertical segment depicts a change in selected cost, not a packet spending time on a vertical path.}
\end{frame}

\begin{frame}[fragile]{Propagation Is Only One Part of Packet Delay}
  \context{A short route through space can still carry packets slowly when other delay terms are large.}
  \begin{columns}[T,onlytextwidth]
    \begin{column}{0.47\textwidth}\raggedright
      \begin{teachingpoints}
        \point{Separate the four common terms.}{\item Propagation moves a signal. Transmission places bits on a link. Queueing waits for service. Processing examines the packet.}
        \point{Use compatible units.}{\item A 1500-byte packet has 12000 bits. On a 10 Mb/s link, transmission takes 1.2 ms before adding other terms.}
        \point{Do not label propagation as RTT.}{\item Round-trip time includes forward and return paths. An application may also wait for servers or additional exchanges.}
      \end{teachingpoints}
    \end{column}
    \begin{column}{0.49\textwidth}\centering
      \figtitle{One illustrative hop: 16.4 ms total}
\begin{tikzpicture}[x=0.90cm,y=0.90cm]
\path[use as bounding box] (-0.25,-0.2) rectangle (7.15,4.5);

\foreach \x/\w/\shade/\lab in {0.15/2.46/Mid/{Propagation},2.61/0.492/Ink/{Tx},3.102/3.69/Soft/{Queue}}{
 \filldraw[fill=\shade,draw=Muted] (\x,2.55) rectangle ++(\w,0.75);
}
\filldraw[fill=Muted,draw=Muted] (6.792,2.55) rectangle ++(0.082,0.75);
\node[lab] at (1.35,3.75) {6 ms};\node[lab] at (2.9,3.75) {1.2};
\node[lab] at (4.8,3.75) {9 ms};
\node[smalllab,anchor=east] at (6.9,1.45) {Processing: 0.2 ms};
\node[smalllab] at (1.35,2.05) {Propagation};
\node[smalllab] at (2.9,1.6) {Transmission};
\node[smalllab] at (4.8,2.05) {Queueing};
\node[lab] at (3.45,0.65) {$D=\sum_{\text{hops}}(d/c+L/C+q+p)$};
\node[smalllab,text=Muted] at (3.45,0.05) {$L$: bits, $C$: bit/s, $q$: wait, $p$: processing};

\end{tikzpicture}
    \end{column}
  \end{columns}
  \assumption{Constructed store-and-forward example. SatNet Edu computes only the propagation term $\sum d/c$.}
  \note{The widths are proportional to 6, 1.2, 9 and 0.2 ms. End-to-end terms depend on the forwarding model. Background for satellite delay components: RFC 2488 section 2. https://www.rfc-editor.org/rfc/rfc2488.html}
\end{frame}

\begin{frame}[fragile]{An Outage Can Leave a Queue After Access Returns}
  \context{Restoring a link stops the outage, but it does not instantly remove accumulated traffic.}
  \begin{columns}[T,onlytextwidth]
    \begin{column}{0.47\textwidth}\raggedright
      \begin{teachingpoints}
        \point{Traffic can accumulate during a gap.}{\item At an arrival rate of 20 Mb/s, a 0.5-second outage builds 10 Mbit if the buffer has enough space.}
        \point{Drain using spare service rate.}{\item After recovery at 50 Mb/s, ongoing arrivals still consume 20 Mb/s. Only 30 Mb/s remains to clear the backlog.}
        \point{Buffer limits change the outcome.}{\item A smaller buffer loses traffic instead. When service is no greater than arrivals, the queue cannot drain.}
      \end{teachingpoints}
    \end{column}
    \begin{column}{0.49\textwidth}\centering
      \figtitle{Recovery continues beyond link restoration}

\begin{tikzpicture}
\begin{axis}[courseaxis,xmin=0,xmax=1.1,ymin=0,ymax=12,
xtick={0,0.5,0.833,1.1},xticklabels={0,0.5,0.833,1.1},
ytick={0,5,10},xlabel={Time (s)},ylabel={Backlog (Mbit)}]
\addplot[Ink,line width=1.1pt] coordinates {(0,0)(0.5,10)(0.833333,0)(1.1,0)};
\draw[Muted,dashed] (axis cs:0.5,0)--(axis cs:0.5,12);
\node[font=\scriptsize] at (axis cs:0.27,11) {Link down};
\node[font=\scriptsize] at (axis cs:0.82,11) {Link restored};
\end{axis}
\end{tikzpicture}
\par\vspace{0.12cm}{\footnotesize Drain time $=10/(50-20)=0.333$ s\par}
    \end{column}
  \end{columns}
  \assumption{Ideal fluid queue, constant arrivals, adequate buffer. This calculation is outside the SatNet Edu model.}
  \note{Queue growth has slope 20 Mbit/s during the outage and minus 30 Mbit/s after recovery. A topology trace alone cannot determine this backlog because arrival and service rates are not modeled.}
\end{frame}

\begin{frame}[fragile]{A Sender Needs Enough Data in Flight}
  \context{A sender waiting for acknowledgements can underuse a fast path when its send window is too small.}
  \begin{columns}[T,onlytextwidth]
    \begin{column}{0.47\textwidth}\raggedright
      \begin{teachingpoints}
        \point{Relate capacity to feedback time.}{\item At bottleneck rate $C$ and round-trip time RTT, the bandwidth-delay product is $C\times\mathrm{RTT}$ bits.}
        \point{A small window can limit the rate.}{\item With at most $W$ unacknowledged bits, an ideal window-limited flow achieves no more than $W/\mathrm{RTT}$.}
        \point{RTT changes the required window.}{\item At 100 Mb/s and 40 ms RTT, 4 Mbit must be in flight to keep the ideal path busy.}
      \end{teachingpoints}
    \end{column}
    \begin{column}{0.49\textwidth}\centering
      \figtitle{100 Mb/s with a 40 ms RTT}
\begin{tikzpicture}[x=0.90cm,y=0.90cm]
\path[use as bounding box] (-0.25,-0.2) rectangle (7.15,4.5);

\node[box] (s) at (0.8,3.55) {Sender};\node[box] (r) at (5.9,3.55) {Receiver};
\draw[flow] (s) to[bend left=20] node[weight,above]{Data in flight}(r);
\draw[flow] (r) to[bend left=20] node[weight,below]{Acknowledgements}(s);
\node[lab] at (3.35,1.75) {$\mathrm{BDP}=100\times0.040=4$ Mbit};
\node[box,text width=5.4cm] at (3.35,0.65) {A 1 Mbit window gives at most\\$1/0.040=25$ Mb/s in this ideal model};

\end{tikzpicture}
    \end{column}
  \end{columns}
  \assumption{Ideal continuous transfer. Congestion, loss, receiver limits, and protocol behavior can reduce throughput further.}
  \note{W is the effective bound on unacknowledged data, not a claim about one particular TCP implementation. 4 Mbit is 500000 bytes. Background: RFC 2488, section 2, bandwidth-delay product. https://www.rfc-editor.org/rfc/rfc2488.html}
\end{frame}

\begin{frame}[fragile]{A Faster New Path Can Reorder Packets}
  \context{Packets sent later can arrive earlier when a flow moves to a shorter-delay path.}
  \begin{columns}[T,onlytextwidth]
    \begin{column}{0.47\textwidth}\raggedright
      \begin{teachingpoints}
        \point{Compare send and arrival times.}{\item Packet 1 leaves at 0 ms on a 30 ms path. Packet 2 leaves at 5 ms on a 10 ms path.}
        \point{Arrival order can reverse.}{\item Packet 2 arrives at 15 ms, before packet 1 at 30 ms. A route change can reorder traffic without losing either packet.}
        \point{Applications depend on timing.}{\item Real-time playback needs timing margin. Reliable transports may buffer out-of-order data before delivery.}
      \end{teachingpoints}
    \end{column}
    \begin{column}{0.49\textwidth}\centering
      \figtitle{Two paths, crossed arrival order}
\begin{tikzpicture}[x=0.90cm,y=0.90cm]
\path[use as bounding box] (-0.25,-0.2) rectangle (7.15,4.5);

\draw[flow] (0.85,4.15)--(0.85,0.15);
\draw[flow] (5.7,4.15)--(5.7,0.15);
\node[lab,above] at (0.85,4.15) {Sender};\node[lab,above] at (5.7,4.15) {Receiver};
\draw[flow,line width=1pt] (0.85,3.65)--node[weight,pos=0.68,below,sloped]{Packet 1: 30 ms}(5.7,0.65);
\draw[flow,dashed,line width=1pt] (0.85,3.15)--node[weight,pos=0.58,above,sloped]{Packet 2: 10 ms}(5.7,2.15);
\node[smalllab,anchor=east] at (0.72,3.65) {0 ms};\node[smalllab,anchor=east] at (0.72,3.15) {5 ms};
\node[smalllab,anchor=west] at (5.85,2.15) {15 ms};\node[smalllab,anchor=west] at (5.85,0.65) {30 ms};
\node[smalllab] at (3.25,0.1) {Time increases downward};

\end{tikzpicture}
    \end{column}
  \end{columns}
  \assumption{Constructed packet-timing example. Serialization is omitted. SatNet Edu does not simulate packet arrivals.}
  \note{Arrival equals send time plus path delay. The newer packet arrives 15 ms earlier. Do not infer packet loss or specific TCP reactions just from the route trace.}
\end{frame}

\begin{frame}[fragile]{A Coarse Recording Can Hide a Short Interruption}
  \context{Recorded snapshots support claims about sample times, not every instant between them.}
  \begin{columns}[T,onlytextwidth]
    \begin{column}{0.47\textwidth}\raggedright
      \begin{teachingpoints}
        \point{State what was sampled.}{\item A 30-second recording can miss an outage entirely when the gap lies between two samples.}
        \point{Refine the observation interval.}{\item Compare 30-second and 5-second recordings using the same scenario. Check whether important transitions become visible.}
        \point{Keep model limits explicit.}{\item SatNet Edu inserts scheduled failure times. Changes caused by geometry still require sufficiently frequent samples.}
      \end{teachingpoints}
    \end{column}
    \begin{column}{0.49\textwidth}\centering
      \figtitle{A gap from 12 to 18 seconds}
\begin{tikzpicture}[x=0.90cm,y=0.90cm]
\path[use as bounding box] (-0.25,-0.2) rectangle (7.15,4.5);

\fill[Soft] (2.9,0.4) rectangle (3.8,4.1);
\node[smalllab] at (3.35,4.35) {6 s gap};
\foreach \yy/\lab in {3.5/{Actual access},2.25/{30 s samples},1/{5 s samples}}{
 \node[smalllab,anchor=east] at (1.05,\yy) {\lab};\draw[faint] (1.1,\yy)--(6.5,\yy);
}
\draw[Ink,line width=3pt] (1.1,3.5)--(2.9,3.5) (3.8,3.5)--(6.5,3.5);
\foreach \x in {1.1,5.6}{\fill[Ink] (\x,2.25) circle (0.06);}
\foreach \x in {1.1,1.85,2.6,4.1,4.85,5.6,6.35}{\fill[Ink] (\x,1) circle (0.06);}
\draw[Ink,fill=Paper] (3.35,1) circle (0.07);
\foreach \x/\t in {1.1/0,2.9/12,3.8/18,5.6/30}{\node[smalllab,below] at (\x,0.18) {\t};}
\node[smalllab] at (3.8,-0.18) {Time (s). Open circle: no access.};

\end{tikzpicture}
    \end{column}
  \end{columns}
  \assumption{Constructed unscheduled gap. A 5-second interval detects this gap but still does not locate its exact start and end.}
  \note{Both 30-second samples at 0 and 30 report access. The 5-second sample at 15 detects an outage. This example concerns changes inferred from geometry, not explicit scheduled failures, whose start and end are inserted into a SatNet Edu trace.}
\end{frame}
\end{document}
